% Band IX, gedruckte S. 112--121.
\begingroup
\clearpage
\noindent112\hfill\textsc{Nachlass.}\hfill\mbox{}\par\bigskip
\begin{center}[6.]\end{center}
Zur Berechnung [von $\psi$ oder $\varphi$] kann auch Folgendes dienen, wenn
\[
\varphi=P+\delta
\]
gesetzt wird:
\[
\begin{aligned}
\psi&=P+\frac{1-ee}{1-ee\sin P^2}\delta+\tfrac34\frac{ee(1-ee)\sin2P}{(1-ee\sin P^2)^2}\delta\delta+\text{ etc.}\\[5pt]
&=\varphi-\frac{ee\cos P^2}{1-ee\sin P^2}\delta+\tfrac34\frac{ee(1-ee)\sin2P}{(1-ee\sin P^2)^2}\delta\delta+\text{ etc.};
\end{aligned}
\]
[wenn]
\[
\psi=P+\varepsilon
\]
gesetzt wird:
\[
\varphi=\psi+\frac{ee\cos P^2}{1-ee\sin P^2}\varepsilon-\tfrac34\frac{ee(1-ee\sin P^2)\sin2P}{(1-ee)^2}\varepsilon\varepsilon-\text{ etc.}
\]
\bigskip
\begin{center}[7.]\end{center}
[Setzt man]
\[
\begin{aligned}
e\sin\varphi&=\sin u\\
e\sin P&=\sin U,
\end{aligned}
\]
[so ist auch]
\[
\psi=\varphi-\cos\tfrac12(\varphi+\psi).Ae\log\frac{\operatorname{tang}(45^\circ+\frac12u)}{\operatorname{tang}(45^\circ+\frac12U)}\ldots,
\]
wo
\[
A=\frac{206\,265}{\mathrm{Mod.}}.
\]
[Für den Abplattungswert $\dfrac1{302{,}68}$ ist] $\log Ae=4{,}586\,3052$.

In sehr guter Näherung [ist]
\[
\varphi=\psi+\frac{ee}{1-ee}(\psi-P)\cos(\tfrac14P+\tfrac34\psi)^2,
\]
genauer
\[
\varphi=\psi+\frac{ee}{1-ee}(\psi-P)\cos(\tfrac14P+\tfrac34\varphi)^2,
\]
noch genauer
\[
\varphi=\psi+\frac{ee(\varphi-P)\sqrt[3]{(\cos\frac12(\varphi-P))}.\cos\frac12(\varphi+P)\cos\frac12(\varphi+\psi)}{1-ee\sin\frac12(P+\varphi)^2},
\]
oder noch genauer, wenn der Nenner $=1-ee\sin P\sin\varphi$ [gesetzt wird].
\clearpage
\noindent\mbox{}\hfill{\small\textsc{Konforme Projektion des Sphäroids auf die Kugel.}}\hfill113\par\bigskip
\begin{center}[8.]\end{center}
Wenn
\[
\begin{aligned}
\operatorname{tang}(45^\circ+\tfrac12\psi)&=\operatorname{tang}(45^\circ+\tfrac12\varphi)\left(\frac{1-e\sin\varphi}{1+e\sin\varphi}\right)^{\frac12e}.\left(\frac{1+e\sin P}{1-e\sin P}\right)^{\frac12e}\\
&=\operatorname{tang}(45^\circ+\tfrac12\varphi).e^u
\end{aligned}
\]
gesetzt wird, so ist
\begingroup\footnotesize
\[
\frac{\mathrm d\psi}{\mathrm du}=\left[\cos\psi=\frac{2\operatorname{tang}(45^\circ+\frac12\psi)}{1+\operatorname{tang}(45^\circ+\frac12\psi)^2}=\right]
\frac{\cos\varphi}{1+\sin\varphi.u+\frac12uu+\frac16\sin\varphi.u^3+\frac1{24}u^4+\frac1{120}\sin\varphi.u^5+\cdots}
\]
\par\endgroup
und daraus
\[
\psi=\varphi+\cos\varphi.u-\tfrac12\cos\varphi\sin\varphi.uu-\tfrac16\cos\varphi\cos2\varphi.u^3\ldots.
\]
[Ferner folgt aus
\[
\mathrm du=-\frac{ee\cos\varphi}{1-ee\sin\varphi^2}\mathrm d\varphi=-ee(1+ee\sin\varphi^2+e^4\sin\varphi^4+\cdots)\mathrm d\sin\varphi:\,]
\]
\[
-u=ee(\sin\varphi-\sin P)+\tfrac13e^4(\sin\varphi^3-\sin P^3)+\tfrac15e^6(\sin\varphi^5-\sin P^5)+\cdots
\]
oder, $\sin\varphi=\sin P+x=\sigma$ gesetzt,
\begingroup\small
\[
\begin{aligned}
u&=-eex-\tfrac13e^4(3x\sigma\sigma-3xx\sigma+x^3)-\tfrac15e^6(5x\sigma^4-10xx\sigma^3+10x^3\sigma\sigma-5x^4\sigma+x^5)\ldots\\
&=-\frac{eex}{1-ee\sigma\sigma}-\ldots.
\end{aligned}
\]
\par\endgroup
[Es ist]
\begingroup\small
\[
\begin{aligned}
\frac{\cos\varphi}{\cos\psi}\bigg[&=\frac{1+\operatorname{tang}(45^\circ+\frac12\psi)^2}{.\,2\operatorname{tang}(45^\circ+\frac12\psi)}\cos\varphi=\cos iu+\sin\varphi\frac{\sin iu}{i}\bigg]\\
&=1+\sin\varphi.u+\tfrac12uu+\tfrac16\sin\varphi.u^3+\tfrac1{24}u^4+\tfrac1{120}\sin\varphi.u^5+\cdots,
\end{aligned}
\]
\[
\log\frac{\cos\varphi}{\cos\psi}=\sin\varphi.u+\tfrac12\cos\varphi^2.uu-\tfrac13\cos\varphi^2\sin\varphi.u^3-\cdots
\]
\par\endgroup
[und]
\begingroup\footnotesize
\[
\sin\psi=\GaussSquareBracketPair{\frac{\operatorname{tang}(45^\circ+\frac12\psi)^2-1}{\operatorname{tang}(45^\circ+\frac12\psi)^2+1}=\frac{\sin\varphi\cos iu+\dfrac{\sin iu}{i}}{\cos iu+\sin\varphi\dfrac{\sin iu}{i}}=}
\frac{\sin\varphi+u+\frac12\sin\varphi.uu+\frac16u^3+\cdots}{1+\sin\varphi.u+\frac12uu+\frac16\sin\varphi.u^3+\cdots},
\]
\[
\sin\varphi-\sin\psi=\GaussSquareBracketPair{-\frac{\cos\varphi^2\dfrac{\sin iu}{i}}{\cos iu+\sin\varphi\dfrac{\sin iu}{i}}=}-\cos\varphi^2.
\frac{u+\frac16u^3+\frac1{120}u^5+\cdots}{1+\sin\varphi.u+\frac12uu+\frac16\sin\varphi.u^3+\cdots},
\]
\[
\frac{\sin\varphi-\sin\psi}{\cos\psi}=\left[-\cos\varphi\frac{\sin iu}{i}=\right]-\cos\varphi.(u+\tfrac16u^3+\tfrac1{120}u^5+\cdots)\left[=\frac{\mathrm d\log n}{\mathrm d\psi}\right].
\]
\par\endgroup
\vfill\noindent{\footnotesize IX.}\hfill15\hspace*{2cm}\par
\clearpage
\noindent114\hfill\textsc{Nachlass.}\hfill\mbox{}\par\bigskip
\begin{center}[9.]\end{center}
[Für $P$ (Göttingen) $=51^\circ31'48''{,}70$ und mit dem Abplattungswert $\dfrac1{302{,}68}$ ergibt sich:]
\begin{center}
\begingroup\renewcommand{\arraystretch}{1.45}\setlength{\tabcolsep}{12pt}
\begin{tabular}[t]{|r|r|r|}
\hline
$\varphi$&$\psi-\varphi$&\\\hline\hline
$48^\circ\ 0'$&$+36''{,}365$&\raisebox{-.7\baselineskip}{\small$1{,}888$}\\
$10$&$34{,}477$&\raisebox{-.7\baselineskip}{\small$1{,}870$}\\
$20$&$32{,}607$&\raisebox{-.7\baselineskip}{\small$1{,}854$}\\
$30$&$30{,}753$&\raisebox{-.7\baselineskip}{\small$1{,}837$}\\
$40$&$28{,}916$&\raisebox{-.7\baselineskip}{\small$1{,}820$}\\
$50$&$27{,}096$&\raisebox{-.7\baselineskip}{\small$1{,}803$}\\
$49\quad0$&$25{,}293$&\raisebox{-.7\baselineskip}{\small$1{,}786$}\\
$10$&$23{,}507$&\raisebox{-.7\baselineskip}{\small$1{,}770$}\\
$20$&$21{,}737$&\raisebox{-.7\baselineskip}{\small$1{,}752$}\\
$30$&$19{,}985$&\raisebox{-.7\baselineskip}{\small$1{,}736$}\\
$40$&$18{,}249$&\raisebox{-.7\baselineskip}{\small$1{,}718$}\\
$50$&$16{,}531$&\raisebox{-.7\baselineskip}{\small$1{,}702$}\\
$50\quad0$&$14{,}829$&\raisebox{-.7\baselineskip}{\small$1{,}684$}\\
$10$&$13{,}145$&\raisebox{-.7\baselineskip}{\small$1{,}668$}\\
$20$&$11{,}477$&\raisebox{-.7\baselineskip}{\small$1{,}650$}\\
$30$&$9{,}827$&\raisebox{-.7\baselineskip}{\small$1{,}634$}\\
$40$&$8{,}193$&\raisebox{-.7\baselineskip}{\small$1{,}616$}\\
$50$&$6{,}577$&\raisebox{-.7\baselineskip}{\small$1{,}600$}\\
$51\quad0$&$4{,}977$&\raisebox{-.7\baselineskip}{\small$1{,}583$}\\
$10$&$3{,}394$&\raisebox{-.7\baselineskip}{\small$1{,}566$}\\
$20$&$1{,}828$&\raisebox{-.7\baselineskip}{\small$1{,}549$}\\
$30$&$+0{,}279$&\raisebox{-.7\baselineskip}{\small$1{,}532$}\\
$40$&$-1{,}253$&\raisebox{-.7\baselineskip}{\small$1{,}516$}\\
$50$&$2{,}769$&\raisebox{-.7\baselineskip}{\small$1{,}499$}\\
$52\quad0$&$4{,}268$&
\end{tabular}\qquad
\begin{tabular}[t]{|r|r|r|}
\hline
$\varphi$&$\psi-\varphi$&\\\hline\hline
$52^\circ\ 0'$&$-4''{,}268$&\raisebox{-.7\baselineskip}{\small$1{,}482$}\\
$10$&$5{,}750$&\raisebox{-.7\baselineskip}{\small$1{,}465$}\\
$20$&$7{,}215$&\raisebox{-.7\baselineskip}{\small$1{,}449$}\\
$30$&$8{,}664$&\raisebox{-.7\baselineskip}{\small$1{,}432$}\\
$40$&$10{,}096$&\raisebox{-.7\baselineskip}{\small$1{,}415$}\\
$50$&$11{,}511$&\raisebox{-.7\baselineskip}{\small$1{,}399$}\\
$53\quad0$&$12{,}910$&\raisebox{-.7\baselineskip}{\small$1{,}382$}\\
$10$&$14{,}292$&\raisebox{-.7\baselineskip}{\small$1{,}364$}\\
$20$&$15{,}656$&\raisebox{-.7\baselineskip}{\small$1{,}348$}\\
$30$&$17{,}004$&\raisebox{-.7\baselineskip}{\small$1{,}331$}\\
$40$&$18{,}335$&\raisebox{-.7\baselineskip}{\small$1{,}315$}\\
$50$&$19{,}650$&\raisebox{-.7\baselineskip}{\small$1{,}298$}\\
$54\quad0$&$20{,}948$&\raisebox{-.7\baselineskip}{\small$1{,}282$}\\
$10$&$22{,}230$&\raisebox{-.7\baselineskip}{\small$1{,}265$}\\
$20$&$23{,}495$&\raisebox{-.7\baselineskip}{\small$1{,}248$}\\
$30$&$24{,}743$&\raisebox{-.7\baselineskip}{\small$1{,}232$}\\
$40$&$25{,}975$&\raisebox{-.7\baselineskip}{\small$1{,}215$}\\
$50$&$27{,}190$&\raisebox{-.7\baselineskip}{\small$1{,}198$}\\
$55\quad0$&$-28{,}388$&
\end{tabular}
\par\endgroup
\end{center}
\clearpage
\noindent\mbox{}\hfill{\small\textsc{Konforme Projektion des Sphäroids auf die Kugel.}}\hfill115\par\bigskip
\begin{center}[10.]\end{center}
Die Korrekturen der berechneten Azimute [auf der Kugel] wegen der Elliptizität sind, [wenn]
\[
\begin{array}{rccl}
\text{beim ersten Punkt:}&(\varphi-\psi)\sin\zeta&=&a\\
\text{in der Mitte:}&\text{»}&=&b\\
\text{beim zweiten Punkt:}&\text{»}&=&c
\end{array}
\]
[und]
\[
s=\text{ganzer Bogen [sowie $\zeta$ dessen südwestliches Azimut ist],}
\]
\[
\begin{array}{rl}
\text{am ersten Punkt:}&(\tfrac13b+\tfrac16a)s\\
\text{»\quad zweiten\quad »\quad :}&-(\tfrac13b+\tfrac16c)s.
\end{array}
\]
Zum \textls[120]{beobachteten} [Azimut] kommt [also] hinzu
\[
\begin{array}{rl}
\text{[am ersten Punkt:]}&-(\tfrac13b+\tfrac16a)s\\
\text{[ »\quad zweiten\quad »\quad :]}&+(\tfrac13b+\tfrac16c)s.
\end{array}
\]
\begin{center}\rule{.18\textwidth}{.4pt}\par\vspace{12mm}\textsc{Bemerkungen.}\end{center}
\begingroup\small
Die vorstehenden Formeln für die konforme Übertragung des Erdellipsoids auf die Kugel und im Wesentlichen auch die unter [II] und [III] folgenden Übertragungsformeln von der Kugel auf die Ebene sind zur selben Zeit entstanden wie die Formeln zur konformen Darstellung des Ellipsoids in der Ebene, die dann bei der hannoverschen Gradmessung und Landesvermessung verwendet wurden. Sie dürften demnach gegen Ende des zweiten Jahrzehnts des vorigen Jahrhunderts entstanden sein. Die Übertragung des Rotationsellipsoids auf die Kugel, wie sie hier gegeben ist und auf Art.~13 der »Allgemeinen Auflösung der Aufgabe, die Theile einer gegebenen Fläche auf einer andern gegebenen Fläche so abzubilden etc.« beruht, wurde später durch die in den »Untersuchungen über Gegenstände der höhern Geodäsie, 1844« mitgeteilte zweite Darstellung abgelöst, die --- wie \textsc{Gauss} selbst sagt (Band IV, S.~262) --- für geodätische Anwendungen noch viel besser geeignet ist.

Die Aufzeichnungen [1], [6], [10] sowie [1], S.~117, bei der stereografischen Projektion und [1], S.~123, bei der Mercatorprojektion folgen in demselben Notizbuch unmittelbar aufeinander.

Die Notizen [2], [3], [7], [8] und [9] finden sich verstreut zwischen Formeln zur konformen Darstellung des Erdellipsoids in der Ebene in einem anderen Notizbuch; [4] und [5] stehen auf einzelnen Blättern zwischen Zahlenrechnungen.
\par\endgroup
\vfill\hfill$15^*$\hspace*{2cm}\par
\clearpage
\noindent116\hfill{\footnotesize\textsc{Bemerkungen.\quad Konforme Projektion des Sphäroids auf die Kugel.}}\par\medskip
\begingroup\small
In [3] ist im Original in der Formel für $\dfrac{\mathrm d\log n}{\mathrm d\psi}$ in dem mit $(\psi-P)^3$ multiplizierten Glied $\tfrac23\cos P^2+\tfrac12\sin P^2$ anstelle von $\tfrac23\cos P^2-\tfrac12\sin P^2$ angegeben; dementsprechend steht auch in der Formel für $\log n$ unter [4] im letzten Glied $\tfrac16\cos P^2+\tfrac18\sin P^2$ anstelle von $\tfrac16\cos P^2-\tfrac18\sin P^2$ und in dem Ausdruck für $\gamma$ unter [5] $\tfrac1{30}\cos P^2+\tfrac1{40}\sin P^2$ anstelle von $\tfrac1{30}\cos P^2-\tfrac1{40}\sin P^2$.

In den Formeln für $\psi$ und $\varphi$ unter [6] hat \textsc{Gauss} in den letzten Gliedern anstelle des Faktors $\tfrac34$ den Faktor $\tfrac14$; bei der Formel $\psi=\varphi-(1)\delta+(2)\delta\delta$ des Beispiels auf S.~119 ist daher (und auch wegen eines kleinen Rechenfehlers) $\log(2)=1{,}86733-10$, wie im Original angegeben, durch $\log(2)=2{,}35885-10$ ersetzt worden.

Die Unterschiede $\psi_1$ und $\psi_2$ zwischen den Azimuten des Bildes der geodätischen Linie des Sphäroids und den Azimuten des größten Kreises auf der Kugel, Art.~[10], ergeben sich nach Band IV, S.~278 u.~f., aus den Gleichungen:
\[
\begin{aligned}
\psi_1&=\tfrac12l_1s+\tfrac16\lambda ss\ldots\\
\psi_2&=-\tfrac12l_1s-\tfrac13\lambda ss\ldots\\
&=-\tfrac12l_2s+\tfrac16\lambda ss\ldots;
\end{aligned}
\]
$l_1$ und $l_2$ sind die Werte von $l=\dfrac{\mathrm dn}{n\,\mathrm d\psi}\sin\zeta$ am Anfangs- und Endpunkt der geodätischen Linie. Ist $l_m$ der Wert von $l$ in der Mitte, so ist
\[
l_m=l_1+\lambda\frac s2\ldots=l_2-\lambda\frac s2\ldots.
\]
Also wird
\[
\psi_1=(\tfrac16l_1+\tfrac13l_m)s\ldots,\qquad\psi_2=-(\tfrac16l_2+\tfrac13l_m)s\ldots.
\]
Nach [1], S.~108, ist $l=\dfrac{\sin\varphi-\sin\psi}{\cos\psi}\sin\zeta=(\varphi-\psi)\sin\zeta-\ldots.$\par
\noindent\hfill\textsc{Krüger, Börsch.}\par
\par\endgroup
\vfill\begin{center}\rule{.24\textwidth}{.4pt}\end{center}
\clearpage
% Unnummerierter Abschnittsbeginn an der Stelle der gedruckten S. 117.
\begin{center}
[II.]\par\bigskip
{\large [STEREOGRAFISCHE PROJEKTION DER KUGEL\\AUF DIE EBENE.]}\par\bigskip
\rule{.1\textwidth}{.4pt}\par\bigskip
[1.]
\end{center}
Wird nun die Kugel durch \textls[110]{stereografische} Projektion auf die Ebene abgebildet, so ergibt sich die Lage des Ortes, dessen Abszisse und Ordinate $x$, $y$ [sind], wie folgt. Man setze
\[
\begin{aligned}
2r\operatorname{tang}(45^\circ-\tfrac12P)+x&=A\cos\alpha&\qquad y&=A\sin\alpha\\
2r\operatorname{tang}(45^\circ+\tfrac12P)-x&=B\cos\beta&y&=B\sin\beta.
\end{aligned}
\]
[$x$ ist nach Süden positiv. $P$ ist die geografische Breite des Zentralpunkts. Sind $n$ und $s$ die Projektionen der beiden Pole und ist $q$ die Projektion eines beliebigen Punkts mit der Breite $\psi$ und der Länge $\lambda$, so gilt $nq=A$ und $sq=B$; $\alpha$ und $\beta$ sind die Winkel, die $A$ und $B$ mit der $x$-Achse bilden.] Dann ist die Länge des Ortes $=\alpha+\beta$ [und zwar wächst die Länge mit $y$]; und die Breite ergibt sich aus
\[
\operatorname{tang}(45^\circ+\tfrac12\psi)=\frac BA\operatorname{tang}(45^\circ-\tfrac12P).
\]
[Ferner ist]
\[
x+y\operatorname{tang}\tfrac12(\alpha-\beta)=2r\operatorname{tang}\tfrac12(P-\psi).
\]
[Die] Neigung des Meridians des Ortes gegen die Abszissenachse [ist] $\alpha-\beta$.
\begin{center}Beispiel.\end{center}
Es sei $P$ [Marseille] $=43^\circ17'50''{,}1$; [die Koordinaten der] Insel Planier [sind]:
\[
x=+5628{,}56\ \text{Tois.}\qquad y=+5679{,}67\ \text{Tois.}
\]
\clearpage
\noindent118\hfill\textsc{Nachlass.}\hfill\mbox{}\par\medskip
[Für die Berechnung von $r$ ist der Abplattungswert $\tfrac1{310}$, vgl. S.~68, verwendet worden.]
\begin{center}
\begingroup\fontsize{8}{12}\selectfont\setlength{\tabcolsep}{3pt}\renewcommand{\arraystretch}{1.4}
\begin{tabular}{@{}r@{ }r@{\quad}r@{ }r@{\quad}l@{}}
$a\ \text{(in Toisen)} .\,.\,.\,.\,.\,.$&$6{,}514\,7609$&$y.\,.\,.\,.\,.\,.\,.\,.\,.\,.$&$3{,}754\,3231$&\\
$\sqrt{(1-ee\sin P^2)}.\,.\,.\,.$&$\underline{9{,}999\,3412}$&$A\cos\alpha.\,.\,.\,.\,.\,.$&$6{,}452\,5266$&$\alpha=6'53''{,}258$\\
$r.\,.\,.\,.\,.\,.\,.\,.\,.\,.\,.\,.\,.\,.$&$6{,}515\,4197$&$B\cos\beta.\,.\,.\,.\,.\,.$&$7{,}181\,0749$&$\underline{\beta=1\ 17{,}210}$\\\cline{3-4}
$2r.\,.\,.\,.\,.\,.\,.\,.\,.\,.\,.\,.\,.\,.$&$6{,}816\,4497$&$\dfrac BA.\,.\,.\,.\,.\,.\,.\,.\,.\,.\,.\,.\,.\,.$&$0{,}728\,5474$&$\lambda=8'10''{,}468$\\
$\operatorname{tang}(45^\circ+\tfrac12P).\,.\,.\,.$&$0{,}364\,7863$&$\operatorname{tang}(45^\circ+\tfrac12P).\,.\,.$&$\underline{0{,}364\,7863}$&\\
$2r\operatorname{tang}(45^\circ-\tfrac12P)=$&$2\,829\,199$&$\operatorname{tang}(45^\circ+\tfrac12\psi).\,.\,.$&$0{,}363\,7611$&\\
$2r\operatorname{tang}(45^\circ+\tfrac12P)=$&$15\,178\,749$&$45^\circ+\tfrac12\psi=$&$66^\circ35'57''{,}71$&\\
$x=$&$5628{,}56$&$\psi=$&$43\ 11\ 55{,}42.$&
\end{tabular}\par\endgroup
\end{center}
Für die Umkehraufgabe dient die Formel
\[
x+iy=2r\frac{\cos\frac12\lambda\sin\frac12(P-\psi)-i\sin\frac12\lambda\cos\frac12(P+\psi)}{\cos\frac12\lambda\cos\frac12(P-\psi)+i\sin\frac12\lambda\sin\frac12(P+\psi)}.
\]
[Dabei ist $\lambda$ nach Osten positiv. Wie vorher ist $x$ nach Süden, $y$ dagegen nach Westen positiv. Ist das nordöstliche Azimut des durch den Anfangspunkt und den Ort $(\psi,\lambda)$ verlaufenden größten Kreises am Anfangspunkt $=T$ und am Ort $=T'$, so ist
\[
\begin{aligned}
\cos\tfrac12(T'+T)\sin\tfrac12\Delta&=\sin\tfrac12(\psi-P)\cos\tfrac12\lambda\\
\sin\tfrac12(T'+T)\sin\tfrac12\Delta&=\cos\tfrac12(\psi+P)\sin\tfrac12\lambda\\
\cos\tfrac12(T'-T)\cos\tfrac12\Delta&=\cos\tfrac12(\psi-P)\cos\tfrac12\lambda\\
\sin\tfrac12(T'-T)\cos\tfrac12\Delta&=\sin\tfrac12(\psi+P)\sin\tfrac12\lambda.
\end{aligned}
\]
Damit wird
\[
x+iy=-2re^{iT}\operatorname{tang}\tfrac12\Delta,
\]
also
\[
x=-2r\operatorname{tang}\frac\Delta2\cos T\qquad y=-2r\operatorname{tang}\frac\Delta2\sin T.\,]
\]
\begin{center}[Beispiel.]\end{center}
\[
\begin{array}{rcl@{\qquad}rl}
\text{Göttingen:}&P=&51^\circ31'55''&\text{Länge}=&27^\circ35'45''\\
\text{Brocken:}&\varphi=&51\ 48\ 12&&\underline{28\ 17\ 1}\\
&&&\lambda=&+41'16''
\end{array}
\]
[Zunächst hat man nach [7], S.~112, und für den Abplattungswert $\tfrac1{310}$]
\clearpage
\noindent\mbox{}\hfill{\small\textsc{Stereografische Projektion der Kugel auf die Ebene.}}\hfill119\par\medskip
\begin{center}
\begingroup\fontsize{8}{11}\selectfont\setlength{\tabcolsep}{3pt}\renewcommand{\arraystretch}{1.4}
\begin{tabular}{@{}r@{ }r@{\quad}r@{ }r@{\quad}l@{}}
$\sin P.\,.\,.$&$9{,}893\,7368$&$\sin\varphi.\,.\,.$&$9{,}895\,3633$&[Man hat auch:]\\
$e.\,.\,.\,.\,.\,.$&$8{,}904\,4836$&$e.\,.\,.\,.\,.\,.$&$8{,}904\,4836$&$\varphi=P+\delta$\\
$U=$&$3^\circ36'9''{,}755$&$u=$&$3^\circ36'58''{,}484$&$\psi=\varphi-(1)\delta+(2)\delta\delta$\\
$\operatorname{tang}(45^\circ+\tfrac12U).\,.\,.$&$0{,}027\,3261$&&&$(1).\,.\,. 7{,}39837-10$\\
$\operatorname{tang}(45^\circ+\tfrac12u).\,.\,.\,.$&$\underline{0{,}027\,4289}$&&&$(2).\,.\,. 2{,}35885-10$\\
&$0{,}000\,1028$&$1028.\,.\,.\,.\,.\,.\,.\,.$&$3{,}01199$&\\
&&$e.\,.\,.\,.\,.\,.\,.\,.\,.\,.\,.\,.$&$8{,}90448$&$[\delta=977''$\\
&&$\dfrac{206\,265}{\mathrm{Mod.}\,10^7}\cos\varphi.\,.\,.$&$\underline{8{,}46788}$&$\psi=\varphi-2''{,}42]$\\
&&$\varphi-\psi.\,.\,.\,.\,.\,.\,.\,.$&$0{,}38435$&\\
&&\multicolumn{2}{c}{$\varphi-\psi=2''{,}42$}&$\psi=51^\circ48'9''{,}58$
\end{tabular}\par\endgroup
\end{center}
\begin{center}
\begingroup\fontsize{8}{11}\selectfont\setlength{\tabcolsep}{3pt}\renewcommand{\arraystretch}{1.4}
\begin{tabular}{@{}r@{ }r@{\quad}r@{ }r@{\quad}l@{}}
$\sin\tfrac12(\psi-P).\,.\,.\,.$&$7{,}37336$&$\cos\tfrac12(\psi+P).\,.\,.\,.$&$9{,}79255$&$\tfrac12(T'+T)=57^\circ36'\ 0''$\\
$\cos\tfrac12\lambda.\,.\,.\,.\,.\,.\,.\,.$&$9{,}99999$&$\sin\tfrac12\lambda.\,.\,.\,.\,.\,.\,.\,.\,.$&$7{,}77829$&$\tfrac12(T'-T)=\phantom{5}0\ 16\ 11$\\
$\cos\tfrac12(\psi-P).\,.\,.$&$0{,}00000$&$\sin\tfrac12(\psi+P).\,.\,.\,.$&$9{,}89455$&$T=57\ 19\ 49$\\\cline{1-2}\cline{3-4}
$\sin\tfrac12\lambda\cos\tfrac12(\psi+P).\,.$&$7{,}57084$&$\sin\tfrac12\lambda\sin\tfrac12(\psi+P).\,.$&$7{,}67284$&\\
$\cos\tfrac12\lambda\sin\tfrac12(\psi-P).\,.$&$\underline{7{,}37335}$&$\cos\mathord{|}\tfrac12\lambda\cos\tfrac12(\psi-P).\,.$&$\underline{9{,}99999}$&\\
$[\operatorname{tang}\tfrac12(T'+T).\,.\,.\,.$&$0{,}19749$&$\operatorname{tang}\tfrac12(T'-T).\,.\,.\,.\,.$&$7{,}67285]$&\\
$\sin\tfrac12(T'+T).\,.\,.\,.\,.\,.$&$9{,}92651$&$\cos\tfrac12(T'-T).\,.\,.\,.\,.$&$0{,}00000$&\\
$\sin\tfrac12\Delta.\,.\,.\,.\,.\,.\,.\,.\,.\,.\,.$&$7{,}64433$&$\cos\tfrac12\Delta.\,.\,.\,.\,.\,.\,.\,.\,.\,.\,.$&$9{,}99999$&
\end{tabular}\par\endgroup
\end{center}
\begin{center}
\begingroup\small\renewcommand{\arraystretch}{1.15}
\begin{tabular}{r@{ }r@{\qquad}l}
$2a.\,.\,.\,.\,.\,.\,.\,.\,.\,.$&$7{,}105\,6109$&\\
$\cos U.\,.\,.\,.\,.\,.\,.$&$\underline{9{,}999\,1408}$&\\
$2r.\,.\,.\,.\,.\,.\,.\,.\,.\,.$&$7{,}106\,4701$&\\
$\operatorname{tang}\tfrac12\Delta.\,.\,.\,.\,.$&$\underline{7{,}64434}$&\\
$2r\operatorname{tang}\tfrac12\Delta.\,.\,.$&$4{,}75081$&$56339\ \mathrm m=\sqrt{(xx+yy)}$\\
$\cos T.\,.\,.\,.\,.\,.\,.$&$9{,}73223$&$-30411\ \text{»}=x$\\
$\sin T.\,.\,.\,.\,.\,.\,.\,.$&$9{,}92521$&$-47427\ \text{»}=y.$
\end{tabular}\par\endgroup
\end{center}
Der doppelte Hilfswinkel des Nenners [nämlich $\tfrac12(T'-T)$] ergibt die Konvergenz der Meridiane oder den obigen Winkel $\alpha-\beta$, hier $32'22''$.

Schließlich ergibt sich die Korrektur des Azimuts des Ortes $x'$, $y'$, vom Ort $x$, $y$ aus gesehen und geradlinig berechnet, aus der Formel:
\[
\operatorname{tang}\mathrm{Corr.}=\frac{x'y-xy'}{4rr+xx'+yy'}\quad\text{genau}\quad=\frac{\dfrac{2rd}{\operatorname{tang}\frac12\delta}\sin\mathrm{Corr.}}{\dfrac{2rd}{\operatorname{tang}\frac12\delta}\cos\mathrm{Corr.}},
\]
\clearpage
\noindent120\hfill\textsc{Nachlass.}\hfill\mbox{}\par\medskip
wenn
\[
\begin{array}{rcccl}
d&\text{die}&\text{Entfernung}&\text{in der Ebene}&\\
\delta&\text{»}&\text{»}&\text{auf der Kugel}&\text{[ist].}
\end{array}
\]
\bigskip
\begin{center}
[2.]\par\medskip
Die stereografische Projektion auf den Horizont eines Ortes,\\dessen Polhöhe $B$ [und dessen] Länge $0$ [ist].
\end{center}
Die Polhöhe des abzubildenden Ortes [sei] $\beta$, [seine] Länge $\lambda$. [Dann ist]
\[
\begin{aligned}
\text{seine Projektion, nördlich:}\quad&\frac{\sin\beta\cos B-\cos\beta\sin B\cos\lambda}{1+\sin B\sin\beta+\cos B\cos\beta\cos\lambda}=y\\[6pt]
\text{östlich:}\quad&\frac{\cos\beta\sin\lambda}{1+\sin B\sin\beta+\cos B\cos\beta\cos\lambda}=x
\end{aligned}
\]
[wenn der Durchmesser der Kugel $=1$ angenommen wird].

Oder
\[
x=\frac{\cos\beta\sin\lambda}{1+\cos\frac12\lambda^2\cos(\beta-B)-\sin\frac12\lambda^2\cos(\beta+B)}
\]
\[
\begin{aligned}
y&=\frac{\cos\frac12\lambda^2\sin(\beta-B)+\sin\frac12\lambda^2\sin(\beta+B)}{1+\cos\frac12\lambda^2\cos(\beta-B)-\sin\frac12\lambda^2\cos(\beta+B)}\\[6pt]
&=\frac{\dfrac1{\cos B}(\sin B+\sin\beta)}{\mathrm{Denomin.}}-\operatorname{tang}B.
\end{aligned}
\]
Der gemeinsame Nenner [ist auch]
\[
=2\cos\tfrac12\lambda^2\cos\tfrac12(\beta-B)^2+2\sin\tfrac12\lambda^2\sin\tfrac12(\beta+B)^2.
\]
1) Die Meridiane werden zu Kreisen mit dem Radius $\dfrac1{\sin\lambda\cos B}$ und dem Mittelpunkt:
\[
\begin{aligned}
y&=-\operatorname{tang}B\\
x&=\frac1{\cos B\operatorname{tang}\lambda}.
\end{aligned}
\]
\bigskip
2) Die Parallelkreise werden zu Kreisen mit dem Radius
\[
\tfrac12\operatorname{cotang}\tfrac12(\beta+B)-\tfrac12\operatorname{tang}\tfrac12(\beta-B)
\]
und dem Mittelpunkt:
\[
\begin{aligned}
y&=\tfrac12\operatorname{cotang}\tfrac12(\beta+B)+\tfrac12\operatorname{tang}\tfrac12(\beta-B)\\
x&=0.
\end{aligned}
\]
\clearpage
\noindent\mbox{}\hfill{\small\textsc{Stereografische Projektion der Kugel auf die Ebene.}}\hfill121\par\bigskip
\begin{center}[3.]\end{center}
[Es sei]
\[
\begin{gathered}
\left.\begin{array}{ll}
2p\, .\, .\, .\, .&\text{Polardistanz}\\
\lambda\, .\, .\, .\, .\, .&\text{Länge}
\end{array}\right\}\ \text{[eines beliebigen Punktes auf der Kugel]}\\[5pt]
2q\, .\, .\, .\, .\ \text{Polardistanz des Zentralpunkts;}
\end{gathered}
\]
[Ferner sei in der Ebene die $x$-Achse das Bild des Anfangsmeridians; ihre positive Richtung weist nach Süden. Die Ordinate $y$ hat dasselbe Vorzeichen wie $\lambda$.

Wenn der Durchmesser der Kugel $=1$ ist, so wird]
\[
x+iy=\frac{\operatorname{tang}p(\cos\lambda+i\sin\lambda)-\operatorname{tang}q}{1+\operatorname{tang}q\operatorname{tang}p(\cos\lambda+i\sin\lambda)}.
\]
[Man setze]
\[
\begin{aligned}
\cos\lambda\operatorname{tang}p-\operatorname{tang}q&=a\cos A\\
\sin\lambda\operatorname{tang}p&=a\sin A\\
1+\operatorname{tang}q\operatorname{tang}p\cos\lambda&=b\cos B\\
\operatorname{tang}q\operatorname{tang}p\sin\lambda&=b\sin B.
\end{aligned}
\]
[Dann wird]
\[
x+iy=\frac ab\{\cos(A-B)+i\sin(A-B)\}.
\]
Noch einfacher
\[
x+iy=\frac{\sin(p-q)\cos\frac12\lambda+i\sin(p+q)\sin\frac12\lambda}{\cos(p-q)\cos\frac12\lambda-i\cos(p+q)\sin\frac12\lambda}.
\]
\begin{center}\rule{.18\textwidth}{.4pt}\par\vspace{10mm}\textsc{Bemerkungen.}\end{center}
\begingroup\small
Wie schon auf S.~115 erwähnt, ist die Notiz [1] in einem Notizbuch enthalten; in den Zahlenbeispielen wurden die Bezeichnungen ergänzt und einige kleine Rechenungenauigkeiten beseitigt. Die Azimutreduktion am Schluss dieser Notiz lässt sich wie folgt herleiten. Es sei $x$ nach Süden positiv, und $y$ habe dasselbe Vorzeichen wie $\lambda$. Dann ist
\[
x=-2r\operatorname{tang}\frac\Delta2\cos T=-\rho\cos T,\qquad y=2r\operatorname{tang}\frac\Delta2\sin T=\rho\sin T,
\]
wobei also $\Delta$ der Großkreisbogen ist, der den abzubildenden Punkt mit dem Zentralpunkt verbindet, und
\par\endgroup
\vfill\noindent{\footnotesize IX.}\hfill16\hspace*{2cm}\par
\endgroup
